\clearpage
\fancyhead[L]{\small INSTRUMENTE.}
\hypertarget{native-429-part-lower}{}
\begin{center}{\Large INSTRUMENTE.}\end{center}
\begin{center}
Astronomisches Jahrbuch für 1822. Seite 235..241. Berlin 1819.
\end{center}

\par\indent Den Kometen habe ich zuerst am 2. Juli in Lauenburg beobachtet, an dem Abend, an dem ihn dort Leutnant \textsc{Zartmann} entdeckte. Ich habe diese Beobachtungen (durch Messung von Abständen zu Sternen mit Spiegelsextanten) jedoch nicht genau reduziert, da ich bald darauf erfuhr, dass der Komet etwa zur gleichen Zeit auch schon andernorts entdeckt und im Meridian beobachtet worden war. Nach meiner Rückkehr habe ich ihn hier bei seiner unteren Kulmination im Meridian beobachtet, zunächst bis zum 26. Juli am Kreis von \textsc{Repsold}. Dort wurden die Durchgänge an den feinen Fäden allerdings allmählich schwer zu beobachten, während die Zenitdistanzen noch mehr Genauigkeit zu erlauben. 
Inzwischen hatte ich das Mittagsfernrohr von \textsc{Reichenbach} (das bei beträchtlich größerer optischer Kraft etwas stärkere Fäden hat) zur Beobachtung bereit gemacht. Daher beobachtete ich die Rektaszensionen vom 27. Juli an mit diesem Instrument, während Herr Professor \textsc{Harding} die Zenitdistanz weiterhin am Kreis von \textsc{Repsold} einstellte. Seit dem 4. August habe ich den Kometen noch an mehreren Abenden (zuletzt am 21. August) im Mittagsfernrohr gesehen. Da die \SourceLeaf{430}Fadendurchgänge aber eher zu erahnen als wirklich zu beobachten waren, lasse ich sie lieber weg, weil sie kein Vertrauen verdienen. Sie bestätigen dennoch die anhaltend gute Übereinstimmung der Elemente von \textsc{Dircksen}.

\Needspace{18\baselineskip}
\begin{center}
\emph{Beobachtungen des Kometen.}
\end{center}

\begin{center}
\begin{tabular}{l|c|c|c}
1819 & Mittlere Zeit & & \\
\hline
Juli 19 & $11^{\mathrm u}\,56^{\mathrm m}\,53^{\mathrm s}$ & $116^\circ\,7'\,26''$ & $+51^\circ\,52'\,12''$ \\
21 & $11\,53\,33$ & $117\,15\,34$ & $51\,54\,58$ \\
22 & $11\,51\,53$ & $117\,49\,34$ & $51\,54\,56$ \\
26 & $11\,44\,14$ & $119\,51\,10$ & $51\,48\,49$ \\
27 & $11\,42\,8$ & $120\,18\,43$ & --- --- --- \\
28 & $11\,40\,2$ & --- --- --- & $51\,43\,46$ \\
30 & $11\,35\,40$ & $121\,38\,52$ & $51\,37\,23$ \\
Aug. 3 & $11\,26\,24$ & $123\,16\,9$ & $51\,23\,25$ \\
4 & $11\,23\,59$ & $123\,38\,55$ & $51\,19\,46$ \\
\end{tabular}
\end{center}

\Needspace{13\baselineskip}
\par\indent Herr \textsc{Dircksen} hat aus diesen Beobachtungen und den ersten Beobachtungen von Herrn Professor \textsc{Struve} in Dorpat die folgenden parabolischen Elemente berechnet.

\begin{center}
\begin{tabular}{ll}
Zeit der Sonnennähe & 1819 Juni 28.\ $17^{\mathrm n}\,49^{\mathrm m}\,9^{\mathrm s}$ Göttinger Zeit \\
Länge der Sonnennähe & . . . . . . . . . . .\ $287^\circ\,6'\,24''8$ \\
Logarithmus des Abstandes & . . . . . . . . . .\ $9.5330800$ \\
Länge des Knotens & . . . . . . . . . . . .\ $273^\circ\,42'\,8''8$ \\
Neigung der Bahn & . . . . . . . . . . . .\ $80\,45\,12.0$ \\
\multicolumn{2}{l}{Bewegung rechtläufig.}
\end{tabular}
\end{center}

\Needspace{17\baselineskip}
\par\indent Der Vergleich der obigen Beobachtungen mit diesen Elementen ergibt folgende Unterschiede:

\begin{center}
\begin{tabular}{l|c|c}
 &  Rektasz.  &  Dekl.  \\
\hline
Juli 19  &  $-15''5$  &  $-3''5$  \\
21  &  $+39.7$  &  $+5.2$  \\
22  &  $-30.0$  &  $+9.5$  \\
26  &  $-37.9$  &  $+26.5$  \\
27  &  $-9.8$  &  --- ---  \\
\hline
  &  Rektasz.  &  Dekl.  \\
\hline
 Juli 28  &  --- ---  &  $+4''0$  \\
 30  &  $-0''2$  &  $-19.2$  \\
 Aug. 3  &  $-9.1$  &  $-0.3$  \\
 4  &  $+3.8$  &  $-2.9$  \\
\end{tabular}
\end{center}

\Needspace{18\baselineskip}
\par\indent Der Vergleich einiger früherer auswärtiger Beobachtungen ergibt folgende Unterschiede:

\begin{center}
\begin{tabular}{l|c|c|l}
Juli 2  &  $-25''1$  &  $-16''6$  &  Bode  \\
3  &  $-1.0$  &  $-1.6$  &  Struve  \\
3  &  $-1.8$  &  $+44.4$  &  Pond  \\
4  &  $-17.3$  &  $+19.6$  &  Struve  \\
4  &  $-48.3$  &  $-65.8$  &  Bode  \\
5  &  $+12.1$  &  $+18.8$  &  Struve  \\
\hline
 Juli 5  &  $-19''1$  &  $-86''7$  &  Pond  \\
 6  &  $-16.6$  &  $+18.0$  &  Struve  \\
 6  &  $-20.8$  &  $-19.9$  &  Bode  \\
 7  &  $-18.1$  &  $+17.3$  &  Struve  \\
 11  &  $+41.1$  &  $-24.4$  &  Bode  \\
 12  &  $-96.3$  &  $-6.3$  &  Bode  \\
\end{tabular}
\end{center}

\par\indent Das sechsfüßige Mittagsfernrohr von \textsc{Reichenbach} ist nun seit Ende Juli täglich in Gebrauch. Eine kurze Nachricht über dieses Instrument wird Ihnen wohl willkommen sein. Das Objektiv von \textsc{Frauenhofer} hat 52 Pariser Linien Öffnungen. Die 4 Okulare vergrößern (nach meiner noch vorläufigen Schätzung) 80-, 100-, 140-, 200-fach. Ich verwende gewöhnlich das stärkste. Dabei ist freilich eine sorgfältige Einstellung nötig, damit man die Sterne nicht verfehlt, denn der Radius des Gesichtsfeldes beträgt nur $4\tfrac{1}{2}$ Minute. Es sind 7 senkrechte Fäden eingezogen, deren Zwischenräume Äquatorsterne innerhalb von $10^{\mathrm s}$ durchlaufen. An diese kurzen Intervalle gewöhnt man sich nach etwas Übung leicht. Man kann in der Zwischenzeit ohne Weiteres die Beobachtungen aufschreiben, das Okular verschieben \SourceLeaf{431}und früh genug fertig sein, um in aller Ruhe die beiden Schläge der Uhr abzuwarten, zwischen die der nächste Fadendurchgang fällt. Man muss nur daran gewöhnt sein, während dieser Handgriffe die Schläge der Uhr sicher weiterzuzählen, wenn man sie nicht von einem Gehilfen zählen lassen kann. Sobald man diese Fertigkeit erworben hat, ist es überaus angenehm, in kurzer Zeit eine größere Zahl von Bestimmungen zu erhalten. Meine eigene Erfahrung zeigt bereits, dass unter dieser Voraussetzung die größere Zahl und der kleinere Abstand der Fäden die Genauigkeit der einzelnen Beobachtungen nicht beeinträchtigen. Aus fast 300 beobachteten Fadendurchgängen von $\gamma$, $\alpha$ und $\beta$ im Adler finde ich für den wahrscheinlichen Fehler \emph{eines} Durchgangs $= 0^{\mathrm{s}}090$, noch etwas weniger als bei meinen Beobachtungen am Kreis von \textsc{Repsold}, wo die 5 Fäden zeitlich über $16^{\mathrm{s}}$ auseinanderliegen. — Das Fernrohr wird mit einem kleinen Kreis eingestellt, der seitlich am Rohr nahe dem Okular befestigt und in Viertelgrade eingeteilt ist. Auf ihm bewegt sich eine Alhidade mit einer Libelle. Der Nonius an dieser Alhidade zeigt einzelne Minuten an. Die Einrichtung ähnelt der am neuen Mittagsfernrohr in Greenwich. Sie bietet den großen Vorteil, dass die Einstellung bei umgelegtem Instrument ebenso leicht ist. Ist die Alhidade einmal eingestellt, kann man auch das Rohr schneller wieder in die richtige Lage bringen, falls es versehentlich ein wenig verrückt wird. Dagegen scheint die Einstellung, da sie \emph{zwei} Handgriffe erfordert, etwas mehr Zeit zu brauchen als bei der gewöhnlichen Einrichtung. Nachts kommt man auch nur schwer ohne einen Gehilfen aus. Herr \textsc{von Reichenbach} wird mir jedoch noch einen zweiten Kreis mit etwas veränderter Einrichtung liefern. Damit entfällt die letzte Unbequemlichkeit, und man gewinnt noch den Vorteil, dass man die beiden Kreise im Voraus auf zwei Sterne einstellen und sie so auch dann bequem beobachten kann, wenn sie sehr schnell aufeinanderfolgen. Ich teile Ihnen noch meine Beobachtungen der letzten Jupiteropposition mit. Dabei habe ich die Rektaszensionen am Mittagsfernrohr beobachtet, während Herr Professor \textsc{Harding} die Zenitdistanzen am Kreis von \textsc{Repsold} einstellte. Diese Beobachtungen können jedoch noch nicht als Maßstab für die mit dem Mittagsfernrohr erreichbare Genauigkeit dienen, weil die damals verwendete Uhr einige Unregelmäßigkeiten im Gang zeigte.

\begin{center}
\begin{tabular}{r|ccc|cc|cc}
1819 & \multicolumn{3}{c|}{Mittlere Zeit} & \multicolumn{2}{c|}{} & \multicolumn{2}{c}{} \\
\hline
Juli 30 & $12^{\mathrm{u}}$ & $31^{\mathrm{m}}$ & $46^{\mathrm{s}}4$ & $315^{\circ}\,42'$ & $47''3$ & $-17^{\circ}\,47'$ & $50''8$ \\
Aug.\ 1 & 12 & 22 & 52.5 & $27$ & 13.2 & $52$ & 34.8 \\
3 & 12 & 13 & 58.7 & $11$ & 41.1 & $57$ & 12.1 \\
4 & 12 & 9 & 31.7 & $3$ & 53.5 & $59$ & 35.9 \\
5 & 12 & 5 & 4.5 & $314\,56$ & 3.1 & $18\,1$ & 49.3 \\
13 & 11 & 29 & 29.1 & $313\,53$ & 50.7 & $19$ & 56.1 \\
\end{tabular}
\end{center}

\par\indent Im vorigen Herbst erhielt ich von Herrn \textsc{Repsold} eine neue, ganz ausgezeichnete Libelle, die unmittelbar auf das Fernrohr gesetzt wird. Der Vergleich der Beobachtungen in entgegengesetzten Lagen des Instruments ergibt den Winkel zwischen der nivellierten Linie und der Sichtlinie. Seitdem habe ich die Beobachtung von etwa 50 Zirkumpolarsternen zu einer meiner Hauptaufgaben gemacht und etwa 1200 Beobachtungen erhalten. Die Ergebnisse für die Polhöhe und die Deklinationen dieser Sterne werde ich bei anderer Gelegenheit bekannt geben. Zuvor müssen aber noch mehrere andere Prüfungen stattfinden. Zum einen werde ich die Teilung mithilfe eines besonderen Apparats, der jetzt an den Pfeilern angebracht ist, aufs Sorgfältigste prüfen. Zum anderen werde ich eine weitere Reihe von Beobachtungen vornehmen, die ich für unerlässlich halte, wenn die \emph{absoluten} Zenitdistanzen zuverlässig sein sollen. Vielfache Erfahrungen aus neuerer Zeit mit Sterndeklinationen und mit der Schiefe der Ekliptik haben gezeigt, dass wir \SourceLeaf{432}über die wahren absoluten Zenitdistanzen noch ganz im Ungewissen sind. Das eine Instrument gibt sie so an, das andere anders. Da wir die Wirkung der \emph{Schwerkraft} auf die Teile des Instruments nicht aufheben können, entstehen je nach Bauart des Instruments durch Biegung und dergleichen \emph{Fehler}, die mit der Zenitdistanz verschieden ausfallen und für die man die Form $a\sin z+b\cos z$ annehmen darf. Bei Instrumenten, die sich umwenden lassen und an denen jeder Stern in beiden Lagen beobachtet wird, ist der zweite Term $b\cos z$ unschädlich. Er verursacht lediglich einen Kollimationsfehler, der je nach Zenitdistanz unterschiedlich ausfällt. Auch bei einem gut ausgeführten Instrument wird ein solcher Term aber immer noch merklich sein. Anders verhält es sich mit dem ersten Term $a\sin z$. Er verfälscht alle Zenitdistanzen, die kleinen wenig, die größeren mehr. In dieser Hinsicht ist jedes Instrument ein Individuum. Zwar kann der Instrumentenbauer den Fehler stark vermindern, doch halte ich es für unmöglich, ihn mit Sicherheit ganz zu beseitigen. Es bleibt nur, ihn durch Beobachtungen zu bestimmen. Der Vergleich der Deklinationen südlicher Sterne, die Herr \textsc{Littrow} aus \textsc{Bessel}s Beobachtungen abgeleitet hat, mit denen von \textsc{Pond} zeigt mir, dass man die Differenz der Werte von $a$ bei den Kreisen von \textsc{Pond} und \textsc{Bessel} ungefähr $=3''$ setzen muss, um die Ergebnisse möglichst gut zur Übereinstimmung zu bringen. Wie groß die einzelnen Werte sind, bleibt unentschieden. Wüsste man \emph{gewiss}, dass der Mittelpunkt der Figur der Sonne einen Großkreis beschreibt und dass die Refraktion für Sonnen- und Sternlicht genau dieselbe ist, so dürfte man den Unterschied ganz oder größtenteils dem Kreis von \textsc{Pond} zuschreiben. Da diese Voraussetzungen jedoch unbewiesen sind, wäre der Schluss eine petitio principii. Ich bemerke nur noch, dass nach dem Vergleich meiner Beobachtungen mit denen von \textsc{Pond} der Wert von $a$ am Kreis von \textsc{Repsold} von dem am Kreis von \textsc{Pond} um $1''5$ abweicht, und zwar auf der dem Kreis von \textsc{Bessel} entgegengesetzten Seite. Dieser Kreis unterscheidet sich daher von dem von \textsc{Repsold} um $4''5$. Darüber braucht man sich keineswegs zu wundern. Nach mehreren Erfahrungen bewirkt am Kreis von \textsc{Repsold} ein Gewicht von 1 Lot, das man beim Objektiv auflegt, eine Biegung, die die Zenitdistanzen am Horizont um $1''$ vermindert. Die Gegengewichte, die \textsc{Repsold} und \textsc{Reichenbach} am Fernrohr angebracht haben, machen diese Biegung zwar sehr klein, bieten aber keine Gewähr, dass sie ganz aufgehoben ist. Das einzige Mittel, den Wert von $a$ zu erfahren, scheint mir eine hinreichend große Zahl von Beobachtungen mithilfe eines Quecksilberhorizonts zu sein. Ich habe bislang erst einige Male Beobachtungen dieser Art am Polarstern angestellt. Sie ergeben zwar im Mittel fast haarscharf $a=0$, doch müssten die Beobachtungen viel zahlreicher sein, damit ich mich dabei beruhigen könnte. Ich kann aber schon glauben, dass der Wert von $a$ am Kreis von \textsc{Repsold} fast $5''$ betragen könnte. Solche Beobachtungen haben übrigens, besonders bei einem großen Instrument, manche praktische Schwierigkeit.

\par\indent Von meinen Beobachtungen am Kreis von \textsc{Repsold} gebe ich hier noch die Rektaszensionen des Polarsterns an, die ich im vorigen Mai bestimmt habe. Der Dezimalbruch des Tages wird von der unteren Kulmination des Datum an gezählt. Der Vergleich ist mit \textsc{Bessel}s \emph{neuen} Tafeln vorgenommen\footnote{Die beiden kleinen Gleichungen bleiben jedoch unberücksichtigt.} (die Unterschiede der Beobachtungen im vorjährigen Jahrbuch bezogen sich auf \textsc{Bessel}s ältere Tafeln). Wegen der Spalte mit der Überschrift Gewicht verweise ich auf das vorjährige Jahrbuch (S. 214 [S. 415 dieses Bandes]).

\Needspace{22\baselineskip}
\SourceLeaf{433}
\begin{center}
\begin{tabular}{r|l|l|l}
1819  &  \shortstack{Scheinb.\\Rektasz.}  &  \shortstack{Verb. Diff. zu\\Bessels Tafeln}  &  Gewicht  \\
\hline
 &  $0^{\mathrm u}\,55^{\mathrm m}$  &   &   \\
Mai 1.\,0  &  $51^{\mathrm s}03$  &  $+0^{\mathrm s}07$  &  5.52  \\
2.75  &  51.27  &  $-$0.40  &  3.43  \\
3.25  &  52.39  &  $+$0.51  &  4.44  \\
6.00  &  52.97  &  $-$0.09  &  3.56  \\
7.25  &  54.33  &  $+$0.72  &  3.75  \\
\hline
 1819  &  \shortstack{Scheinb.\\Rektasz.}  &  \shortstack{Verb. Diff. zu\\Bessels Tafeln}  &  Gewicht  \\
\hline
  &  $0^{\mathrm u}\,55^{\mathrm m}$  &   &   \\
 Mai 8.25  &  $54^{\mathrm s}83$  &  $+0^{\mathrm s}76$  &  3.43  \\
 9.75  &  54.95  &  $+$0.17  &  3.00  \\
 10.75  &  53.29  &  $-$1.98  &  3.40  \\
 16.50  &  57.21  &  $-$1.06  &  3.00  \\
 21.25  &  60.96  &  $+$0.03  &  2.67  \\
\end{tabular}
\end{center}

Im Mittel wäre demnach die Korrektur nur $= -0^{\mathrm s}03$.

\par\indent Seit Kurzem ist nun auch der Meridiankreis von \textsc{Reichenbach} angekommen und kann bald aufgestellt werden. Zum Mittagsfernrohr bemerke ich noch, dass seine optische Kraft außerordentlich groß ist. Den kleinen Stern, der $5'$ vom Pol entfernt ist, sehe ich in der Dämmerung, wenn noch keine Fadenbeleuchtung nötig ist. Das Umlegen geschieht in der Luft mit einem Apparat, den der hiesige Herr Inspektor \textsc{Rumpf} sehr schön ausgeführt hat. Ein Seil läuft oben über eine Rolle, von dort seitwärts zu einer zweiten Rolle und dann nach unten auf die Welle einer Winde mit Sperre und dergleichen, auf die es aufgewickelt ist. Damit hebt man das Instrument nur etwa 1 Zoll an. Darauf dreht man es frei schwebend um und lässt es in der entgegengesetzten Lage wieder ein. Nach Gebrauch wird die Rolle durch bloßes Ziehen an einer Schnur zur Seite geschoben, damit sie bei Zenitbeobachtungen nicht stört. Der ganze Vorgang dauert etwa 5--6 Minuten, sodass ich beim Umlegen während der Kulmination des Polarsterns von 7 Fäden nur einen verliere. Sobald ich den zweiten Einstellkreis erhalten habe, lässt sich der Vorgang leicht so schnell ausführen, dass gar kein Faden verloren geht. Beim Kreis von \textsc{Repsold} dauert das Umlegen immer über eine halbe Stunde. Beim Kreis von \textsc{Reichenbach} wird es auf andere Weise mit einer besonderen Maschine geschehen, und der Instrumentenbauer glaubt, dass jedes Mal 8 Minuten genügen werden.
